\chapter{METODOLOGI}
\label{sec:chap3_metodologi}

\section{Metodologi Penelitian}

Metodologi penelitian yang digunakan dalam penelitian ini terdiri dari beberapa tahapan, yaitu analisis kebutuhan sistem, perancangan sistem, implementasi sistem, serta pengujian sistem. Tahapan-tahapan tersebut dilakukan untuk mengembangkan dan meningkatkan platform monitoring CCTV berbasis ESP32-CAM (MIOTVision) yang mampu melakukan pemantauan secara real-time melalui integrasi website dan aplikasi mobile yang lebih stabil.

\section{Analisis Kebutuhan Sistem}

Pada tahap ini dilakukan identifikasi kebutuhan sistem yang diperlukan untuk mendukung proses monitoring CCTV berbasis Internet of Things (IoT). Kebutuhan sistem meliputi kebutuhan perangkat keras, perangkat lunak, serta kebutuhan fungsional yang harus dipenuhi oleh sistem yang dikembangkan.

\subsection{Kebutuhan Perangkat Keras}

Perangkat keras yang digunakan pada penelitian ini ditunjukkan pada Tabel \ref{tab:hardware}.

\begin{table}[H]
	\centering
	\caption{Kebutuhan Perangkat Keras}
	\label{tab:hardware}
	\begin{tabular}{|c|l|l|}
		\hline
		No & Perangkat & Fungsi \\
		\hline
		1 & ESP32-CAM & Pengambilan citra \\
		2 & Intel NUC & Server utama sistem \\
		3 & Router/Wi-Fi & Koneksi jaringan \\
		4 & Smartphone & Monitoring dan notifikasi \\
		5 & Laptop & Pengembangan sistem \\
		\hline
	\end{tabular}
\end{table}

\subsection{Kebutuhan Perangkat Lunak}

Perangkat lunak yang digunakan pada penelitian ini ditunjukkan pada Tabel \ref{tab:software}.

\begin{table}[H]
	\centering
	\caption{Kebutuhan Perangkat Lunak}
	\label{tab:software}
	\begin{tabular}{|c|l|}
		\hline
		No & Perangkat Lunak \\
		\hline
		1 & Debian 12 \\
		2 & Docker \\
		3 & Laravel \\
		4 & Flutter \\
		5 & MySQL \\
		6 & MinIO \\
		7 & EMQX \\
		8 & Webmin \\
		9 & Cloudflared \\
		\hline
	\end{tabular}
\end{table}

\section{Perancangan Sistem}

Tahap perancangan dilakukan untuk menentukan arsitektur sistem, alur komunikasi data, serta hubungan antar komponen yang digunakan dalam penelitian.

\subsection{Arsitektur Sistem}

Arsitektur sistem yang dibangun memperluas infrastruktur MIOTVision terdahulu, yang mengintegrasikan perangkat ESP32-CAM, broker MQTT (EMQX), server backend Laravel, basis data MySQL, object storage MinIO, serta antarmuka website dan aplikasi mobile secara terpusat.

%\begin{figure}[H]
%	\centering
%	\includegraphics[width=0.9\textwidth]{gambar/arsitektur_sistem.pdf}
%	\caption{Arsitektur Sistem}
%	\label{fig:arsitektur}
%%\end{figure}

Pada arsitektur tersebut, perangkat ESP32-CAM melakukan akuisisi citra, pengumpulan telemetri, dan publikasi data melalui protokol MQTT over WSS menuju broker EMQX. Server backend Laravel yang bertindak sebagai subscriber menerima data citra tersebut lalu menyimpannya ke object storage MinIO secara terpusat. Selanjutnya, backend Laravel meneruskan permintaan analisis secara asinkron ke FastAPI Detection Service yang berjalan secara terisolasi. Di dalam FastAPI Detection Service, QueueWorker memproses antrean citra menggunakan BackgroundTasks secara asinkron. Citra diproses melalui modul deteksi gerakan (motion detection) berbasis BackgroundSubtractorMOG2 dengan motion thresholding. Jika terdeteksi adanya gerakan, sistem memicu event deteksi gerakan (motion\_events) dan melanjutkan citra ke modul deteksi penyusup (intruder detection) menggunakan model YOLO11n. Apabila objek manusia teridentifikasi, event deteksi penyusup (detection\_events) akan dikirimkan kembali ke Laravel untuk dipersisten ke MySQL, disiarkan secara real-time ke web browser klien via Laravel Reverb, serta dikirimkan sebagai push notification ke aplikasi mobile Flutter.

\subsection{Perancangan Alur Komunikasi Data}

Proses komunikasi data berjalan sesuai alur deteksi terintegrasi: ESP32-CAM melakukan akuisisi citra and publikasi MQTT over WSS menuju broker EMQX. Pesan tersebut diterima oleh backend Laravel untuk disimpan ke object storage MinIO. Citra pada MinIO kemudian diakses oleh FastAPI Detection Service secara asinkron. Di dalam service tersebut, modul motion detection berbasis BackgroundSubtractorMOG2 mengevaluasi perubahan piksel. Jika ambang batas gerakan terlampaui, diproduksi motion\_events dan citra diteruskan ke model YOLO11n untuk klasifikasi manusia. Jika terklasifikasi sebagai penyusup, dihasilkan detection\_events yang dikirimkan ke Laravel API persistence, disimpan ke MySQL database, disiarkan via Laravel Reverb ke website monitoring, dan dikirimkan sebagai push notification ke aplikasi mobile Flutter. Selain itu, alur komunikasi juga mencakup topik MQTT \texttt{ws/camera/\{uuid\}/config} untuk pertukaran status konfigurasi perangkat secara dua arah dan \texttt{ws/camera/\{uuid\}/ota} untuk instruksi pembaruan firmware secara asinkron.

\subsection{Perancangan Basis Data}

Basis data dirancang untuk menyimpan informasi pengguna, perangkat kamera, data citra, data telemetri, dan data notifikasi.

Tabel utama yang digunakan meliputi:

\begin{enumerate}
	\item Tabel \texttt{users}
	\item Tabel \texttt{cameras} (menyimpan parameter sinkronisasi: \texttt{desired\_config}, \texttt{current\_config}, \texttt{desired\_config\_hash}, \texttt{current\_config\_hash}, dan \texttt{pending\_changes})
	\item Tabel \texttt{image\_records}
	\item Tabel \texttt{camera\_telemetry}
	\item Tabel \texttt{notifications}
	\item Tabel \texttt{ota\_firmwares} (metadata biner firmware OTA)
	\item Tabel \texttt{ota\_deployments} (riwayat deployment pembaruan)
\end{enumerate}

\subsection{Perancangan Sinkronisasi Konfigurasi Kamera}

Perancangan sinkronisasi konfigurasi dirancang agar perubahan parameter kamera pada dashboard web dapat diterapkan ke perangkat ESP32-CAM secara asinkron. Perubahan parameter disimpan pada kolom \texttt{desired\_config} dan hash-nya pada \texttt{desired\_config\_hash}, lalu dikirim melalui topik MQTT \texttt{ws/camera/\{uuid\}/config}. Mikrokontroler menerapkan pengaturan tersebut dan mengirim balik hash aktualnya (\texttt{current\_config\_hash}) untuk mematikan status \texttt{pending\_changes}.

\subsection{Perancangan Motion Detection}

Modul motion detection berjalan secara asinkron dalam FastAPI Detection Service di server Intel NUC. Modul ini mengambil gambar dari MinIO dan memprosesnya menggunakan BackgroundSubtractorMOG2 untuk segmentasi latar belakang. Apabila luas area gerakan melebihi parameter motion thresholding yang ditentukan, modul ini mengonfirmasi adanya pergerakan (motion\_events) dan secara otomatis meneruskan frame citra tersebut ke model YOLO11n.

\subsection{Perancangan Monitoring Telemetri}

Monitoring telemetri digunakan untuk memantau kondisi operasional perangkat ESP32-CAM. Informasi yang dikirimkan secara berkala meliputi nilai RSSI, uptime perangkat, free heap, reset reason, capture success rate, dan publish success rate. Data tersebut digunakan untuk membantu proses pemeliharaan dan identifikasi gangguan perangkat.

\section{Implementasi Sistem}

Tahap implementasi dilakukan dengan merealisasikan seluruh rancangan yang telah dibuat ke dalam bentuk sistem yang dapat dijalankan.

\subsection{Implementasi Infrastruktur Server}

Server utama menggunakan Intel NUC dengan sistem operasi Debian 12. Seluruh layanan dijalankan menggunakan Docker Container untuk mempermudah proses deployment dan pengelolaan sistem.

Layanan kontainer Docker yang dijalankan meliputi:

\begin{enumerate}
	\item \texttt{cctv-app} (Backend Laravel)
	\item \texttt{cctv-web} (Web Server Nginx)
	\item \texttt{cctv-db} (MySQL Database)
	\item \texttt{cctv-emqx} (EMQX MQTT Broker)
	\item \texttt{cctv-minio} (MinIO Object Storage)
	\item \texttt{cctv-reverb} (Laravel Reverb WebSocket)
	\item \texttt{cctv-detection-service} (FastAPI Detection Service)
	\item \texttt{cctv-phpmyadmin} (Database Client phpMyAdmin)
\end{enumerate}

\begin{figure}[H]
\centering
[PLACEHOLDER IMAGE]
\caption{Skema Infrastruktur Kontainer Server}
\label{fig:docker_deployment}
\end{figure}

\subsection{Implementasi Cloudflared Tunnel}

Cloudflared Tunnel digunakan untuk menyediakan akses publik terhadap website dan API tanpa membuka port secara langsung pada router. Konfigurasi dilakukan dengan memetakan lalu lintas internet aman (HTTPS) ke masing-masing kontainer lokal menggunakan domain publik: \texttt{cctv.miot-its.org} (aplikasi utama), \texttt{broker.miot-its.org} (WSS broker), \texttt{minio.miot-its.org} (object storage), \texttt{pma.miot-its.org} (phpMyAdmin), dan \texttt{web.miot-its.org} (monitoring frontend).

\begin{figure}[H]
\centering
[PLACEHOLDER IMAGE]
\caption{Diagram Alur Akses Aman Cloudflared Tunnel}
\label{fig:cloudflared_tunnel}
\end{figure}

\subsection{Implementasi Webmin}

Webmin digunakan untuk membantu proses administrasi server melalui antarmuka berbasis web. Fitur ini digunakan untuk melakukan pemantauan kondisi server, pengelolaan layanan, serta konfigurasi sistem operasi.

\begin{figure}[H]
\centering
[PLACEHOLDER IMAGE]
\caption{Dashboard Administrasi Webmin}
\label{fig:webmin_dashboard}
\end{figure}

\subsection{Implementasi Website Monitoring}

Website monitoring dikembangkan menggunakan framework frontend yang terhubung dengan backend melalui WebSocket. Website mampu menampilkan informasi perangkat, status kamera, data telemetri, gambar hasil tangkapan kamera, serta notifikasi aktivitas secara real-time.

\begin{figure}[H]
\centering
[PLACEHOLDER IMAGE]
\caption{Halaman Dashboard Website MIOTVision}
\label{fig:website_monitoring}
\end{figure}

\subsection{Implementasi Aplikasi Mobile}

Aplikasi mobile berbasis Flutter yang telah dirintis sebelumnya ditingkatkan fungsionalitasnya untuk mendukung penerimaan push notification secara real-time, monitoring telemetri perangkat secara visual, serta integrasi status operasional CCTV.

\begin{figure}[H]
\centering
[PLACEHOLDER IMAGE]
\caption{Antarmuka Aplikasi Mobile Flutter}
\label{fig:flutter_app}
\end{figure}

\subsection{Implementasi Deteksi Gerakan dan Deteksi Penyusup}

Layanan deteksi diimplementasikan pada kontainer terpisah menggunakan framework FastAPI dengan antrean asinkron pada \texttt{worker.py} dan fungsi visi komputer pada \texttt{handlers/detection.py} berbasis \texttt{BackgroundTasks}. Deteksi gerakan dilakukan oleh algoritma BackgroundSubtractorMOG2 dengan konfigurasi parameter \texttt{history=500}, \texttt{varThreshold=16}, dan \texttt{detectShadows=False} untuk memicu \texttt{motion\_events}. Citra kemudian ditarik dari bucket \texttt{cctv/} di bawah direktori \texttt{camera/\{uuid\}} pada MinIO untuk diproses oleh model YOLO11n. Klasifikasi penyusup dilakukan dengan memfilter objek manusia menggunakan kondisi \texttt{class\_id == 0}. Ketika terverifikasi, sistem menghasilkan \texttt{detection\_events} yang dikirim ke Laravel API persistence untuk disimpan di MySQL.

\begin{figure}[H]
\centering
[PLACEHOLDER IMAGE]
\caption{Diagram Alur Deteksi Asinkron Detection Service}
\label{fig:detection_service}
\end{figure}

\subsection{Implementasi Pemantauan Telemetri}

Pemantauan telemetri perangkat dilakukan secara periodik dengan mengirimkan data operasional ESP32-CAM ke server melalui protokol MQTT. Data yang dikirimkan meliputi parameter Received Signal Strength Indicator (RSSI), uptime perangkat, free heap memory, reset reason, serta tingkat keberhasilan pengambilan gambar dan publikasi data.

\begin{figure}[H]
\centering
[PLACEHOLDER IMAGE]
\caption{Visualisasi Widget Telemetri Perangkat}
\label{fig:telemetry_widget}
\end{figure}

\subsection{Implementasi Notifikasi Real-time}

Notifikasi real-time diimplementasikan menggunakan layanan push notification untuk mengirimkan peringatan keamanan instan ke perangkat mobile pengguna segera setelah terdeteksi adanya penyusupan atau gerakan mencurigakan pada area pengawasan.

\begin{figure}[H]
\centering
[PLACEHOLDER IMAGE]
\caption{Penerimaan Push Notification pada Handphone}
\label{fig:push_notification}
\end{figure}

\subsection{Integrasi Laravel Reverb}

Laravel Reverb diintegrasikan pada backend untuk menangani komunikasi WebSocket berskala tinggi secara asinkron. Hal ini memungkinkan server mengirimkan pembaruan status perangkat, data telemetri terbaru, dan notifikasi deteksi langsung ke dashboard website secara real-time tanpa overhead HTTP polling.

\subsection{Stabilisasi MQTT over WSS}

Stabilisasi komunikasi MQTT over WebSocket Secure (WSS) dilakukan untuk menjamin transmisi data yang aman dan andal antara perangkat ESP32-CAM and broker EMQX melalui port web standar dengan dukungan mekanisme auto-reconnect.

\subsection{Implementasi Pembaruan Firmware OTA}

Dukungan pembaruan firmware secara nirkabel (Over-the-Air/OTA) diimplementasikan agar perangkat ESP32-CAM dapat mengunduh berkas biner firmware baru secara asinkron dari server backend Laravel dan melakukan pembaruan kode program secara otomatis.

\begin{figure}[H]
\centering
[PLACEHOLDER IMAGE]
\caption{Halaman Manajemen Firmware OTA pada Backend}
\label{fig:ota_management}
\end{figure}

\subsection{Implementasi Kebijakan Retensi MySQL}

Mekaniseme retensi basis data MySQL diimplementasikan dengan membuat prosedur terjadwal (event scheduler) untuk membersihkan data log telemetri dan riwayat deteksi yang telah melewati batas waktu penyimpanan tertentu agar ukuran basis data tetap optimal.

\subsection{Implementasi Pembersihan Otomatis MinIO}

Pembersihan otomatis pada MinIO Object Storage diimplementasikan menggunakan perintah konsol Laravel \texttt{cleanup-image-records} yang dijalankan secara berkala oleh Laravel Scheduler untuk menghapus berkas citra deteksi lama secara terjadwal di bawah direktori \texttt{camera/\{uuid\}} pada bucket \texttt{cctv/} guna menghemat ruang penyimpanan. Selain itu, Laravel Scheduler juga menjalankan Artisan Command \texttt{process-scheduled-ota-deployments} secara berkala untuk mengeksekusi antrean pembaruan firmware kamera dari direktori \texttt{firmware/\{version\}}.

\subsection{Deployment Layanan Deteksi}

Layanan deteksi di-deploy sebagai kontainer Docker terpisah (FastAPI Detection Service) pada server Intel NUC. Konfigurasi container menggunakan Dockerfile dan Docker Compose untuk memisahkan alur kerja visi komputer (BackgroundSubtractorMOG2 dan YOLO11n) dari server Laravel backend, guna memastikan isolasi resource CPU/RAM.

\section{Pengujian Sistem}

Pengujian dilakukan untuk memastikan seluruh fungsi sistem berjalan sesuai dengan kebutuhan yang telah ditentukan.

Pengujian yang dilakukan meliputi:

\begin{enumerate}
	\item Pengujian koneksi perangkat ESP32-CAM.
	\item Pengujian komunikasi MQTT.
	\item Pengujian penyimpanan data pada MySQL dan MinIO.
	\item Pengujian monitoring real-time menggunakan WebSocket.
	\item Pengujian motion detection.
	\item Pengujian monitoring telemetri perangkat.
	\item Pengujian push notification pada aplikasi mobile.
	\item Pengujian akses publik menggunakan Cloudflared Tunnel.
\end{enumerate}